本章确定后续散射问题所需的背景、传播相位与谱坐标. 参数始终满足 $Q>0$、$a\ne0$、$\eta_0\ne0$, 且 $\alpha,\beta$ 为实数.

\section{平面波背景与介质约束}
\begin{proposition}[相容平面波与 Bloch 长度]\label{prop:bg}\label{prop:bloch}
系统 \eqref{eq:HMB} 具有平面波解
\begin{equation}\label{eq:background}
q_0=Qe^{-\ii\Omega z},\qquad
p_0=-\dfrac{\ii\eta_0}{a}q_0,\qquad
\eta=\eta_0,\qquad
\Omega=-\alpha Q^2+\dfrac{2\eta_0}{a}.
\end{equation}
任一满足介质方程的经典解均有
\begin{equation}\label{eq:bloch}
\partial_x(|p|^2+\eta^2)=0.
\end{equation}
若其端态为 \eqref{eq:targetboundary}, 则
\begin{equation}\label{eq:C}
|p|^2+\eta^2=\eta_0^2\left(1+\dfrac{Q^2}{a^2}\right).
\end{equation}
\end{proposition}
\begin{proof}
对不依赖 $x$ 的场, 第二条方程给出 $p_0=-\dfrac{\ii\eta_0q_0}{a}$, 第三条方程随之成立. 第一条方程化为
\[
(q_0)_z=\ii\left(\alpha Q^2-\dfrac{2\eta_0}{a}\right)q_0=-\ii\Omega q_0,
\]
得到所述平面波. 对一般经典解, 利用 $a,\eta$ 为实数,
\begin{align*}
\partial_x(|p|^2+\eta^2)
&=(2\eta q-2\ii ap)\conj p
 +p(2\eta\conj q+2\ii a\conj p)
 -2\eta(\conj q p+q\conj p)=0.
\end{align*}
故该量仅依赖 $z$; 在任一空间端点取极限即得 \eqref{eq:C}.
\end{proof}

\section{Lax 表示与零曲率条件}
采用文献\cite{LHP2013}式 (2.2)--(2.4) 的 Lax 对, 变量转换为 $t_{\rm source}=x$、$E=q$、$\omega=-a$, 其余参数不缩放.
令
\[
A(q)=\begin{pmatrix}0&q\\-\conj q&0\end{pmatrix},\quad
R=\begin{pmatrix}\eta&-p\\-\conj p&-\eta\end{pmatrix},\quad
\sg=\begin{pmatrix}1&0\\0&-1\end{pmatrix}.
\]
相应矩阵为
\begin{equation}\label{eq:lax}
\Phi_x=U\Phi,\quad U=-\ii\lambda\sg+A(q),\qquad
\Phi_z=V\Phi,\quad V=\sum\limits_{j=0}^{3}\lambda^jV_j+
\dfrac{\ii R}{\lambda-a}.
\end{equation}
全部系数为
\begin{align}
V_3&=4\ii\beta\sg,\qquad
V_2=-\ii\alpha\sg-4\beta A(q),\label{eq:V32}\\
V_1&=\begin{pmatrix}
-2\ii\beta|q|^2&\alpha q-2\ii\beta q_x\\
-\alpha\conj q-2\ii\beta\conj{q_x}&2\ii\beta|q|^2
\end{pmatrix},\label{eq:V1}\\
V_0&=\begin{pmatrix}d&b\\c&-d\end{pmatrix},\label{eq:V0}\\
d&=\dfrac{\ii\alpha}{2}|q|^2-\beta(q\conj{q_x}-q_x\conj q),\nonumber\\
b&=2\beta|q|^2q+\dfrac{\ii\alpha}{2}q_x+\beta q_{xx},\nonumber\\
c&=-2\beta|q|^2\conj q+\dfrac{\ii\alpha}{2}\conj{q_x}
-\beta\conj{q_{xx}}.\nonumber
\end{align}

\begin{proposition}[Lax 相容性]\label{prop:zc}
对足够光滑的 $q,p$ 和实值 $\eta$, 式 \eqref{eq:lax} 的零曲率条件
\begin{equation}\label{eq:zc}
U_z-V_x+[U,V]=0
\end{equation}
等价于 \eqref{eq:HMB} 及其复共轭方程. 其矩阵满足约化
\begin{equation}\label{eq:antihermitian}
U(\conj\lambda)^\dagger=-U(\lambda),\qquad
V(\conj\lambda)^\dagger=-V(\lambda).
\end{equation}
\end{proposition}
\begin{proof}
各多项式系数满足相应的反 Hermitian 约化; $R^\dagger=R$ 则给出有理项的约化. 以下比较零曲率表达式的极点与多项式系数.
写 $V_{\rm poly}=\displaystyle\sum\limits_{j=0}^3\lambda^jV_j$. 由
$U(\lambda)=U(a)-\ii(\lambda-a)\sg$, 有理部分为
\begin{equation}\label{eq:rationalzc}
-\dfrac{\ii R_x}{\lambda-a}
+\dfrac{\ii[U(\lambda),R]}{\lambda-a}
=\dfrac{\ii([U(a),R]-R_x)}{\lambda-a}+[\sg,R].
\end{equation}
矩阵相乘给出
\[
[U(a),R]=
\begin{pmatrix}
-(q\conj p+\conj q p)&2\ii ap-2\eta q\\
-2\ii a\conj p-2\eta\conj q&q\conj p+\conj q p
\end{pmatrix}.
\]
极点留数为零恰好给出两条介质方程及其共轭.

按 $\lambda^4,\lambda^3,\lambda^2,\lambda$ 比较多项式部分, 得到
\begin{align*}
&[-\ii\sg,V_3]=0,\\
&[-\ii\sg,V_2]+[A,V_3]=0,\\
&-(V_2)_x+[-\ii\sg,V_1]+[A,V_2]=0,\\
&-(V_1)_x+[-\ii\sg,V_0]+[A,V_1]=0.
\end{align*}
将上述系数代入, 四式恒等成立. 例如二次项的 $(1,2)$ 元为
$4\beta q_x-2\ii(\alpha q-2\ii\beta q_x)+2\ii\alpha q=0$.
一次项的两个独立分量分别为: 其 $(1,1)$ 元为
\[
2\ii\beta(|q|^2)_x
+q(-\alpha\conj q-2\ii\beta\conj{q_x})
+(\alpha q-2\ii\beta q_x)\conj q=0,
\]
而 $(1,2)$ 元为
\begin{align*}
&-\alpha q_x+2\ii\beta q_{xx}-2\ii b+4\ii\beta|q|^2q\\
&\quad=-\alpha q_x+2\ii\beta q_{xx}
-2\ii\left(2\beta|q|^2q+\dfrac{\ii\alpha}{2}q_x+\beta q_{xx}\right)
+4\ii\beta|q|^2q=0.
\end{align*}
四次项是对角矩阵的交换子, 故为零; 三次项为
$4\ii\beta[\sg,A]+4\ii\beta[A,\sg]=0$. 其余分量由迹零与反 Hermitian 约化确定.
常数项 $(1,2)$ 元为 $q_z-b_x-2dq-2p$. 展开
\begin{align*}
b_x+2dq
&=\beta q_{xxx}+\dfrac{\ii\alpha}{2}q_{xx}
 +2\beta(|q|^2q)_x+\ii\alpha|q|^2q
 -2\beta q(q\conj{q_x}-q_x\conj q)\\
&=\beta q_{xxx}+\dfrac{\ii\alpha}{2}q_{xx}
 +6\beta|q|^2q_x+\ii\alpha|q|^2q.
\end{align*}
这给出光场方程. 对常数对角元, 展开得
\begin{align*}
qc+b\conj q
&=\dfrac{\ii\alpha}{2}(q\conj{q_x}+q_x\conj q)
+\beta(q_{xx}\conj q-q\conj{q_{xx}})\\
&=\partial_x\left(\dfrac{\ii\alpha}{2}|q|^2
-\beta(q\conj{q_x}-q_x\conj q)\right)=d_x.
\end{align*}
因此该分量恒等为零. $(2,2)$ 元是其负值, $(2,1)$ 元由反 Hermitian 约化给出光场方程的负共轭. 因此零曲率条件与式 \eqref{eq:HMB} 等价.
\end{proof}

\section{常背景规范与传播相位}
通过去除共同平面波相位, 将两端归一化为同一个常矩阵. 令
\begin{equation}\label{eq:gauge}
G(z)=e^{-\frac{\ii\Omega z}{2}\sg},\quad
\Phi=G\Psi,\quad u=e^{\ii\Omega z}q,\quad v=e^{\ii\Omega z}p.
\end{equation}
以下 $R$ 专指由去相位变量构成的矩阵
$R=\left(\begin{smallmatrix}\eta&-v\\-\conj v&-\eta\end{smallmatrix}\right)$, 并记
$U_\infty=-\ii\lambda\sg+A(Q)$.

\begin{proposition}[背景规范变换]\label{prop:gauge}
变换 \eqref{eq:gauge} 将 Lax 对化为
\begin{equation}\label{eq:gaugelax}
\Psi_x=\widehat U\Psi,\quad
\Psi_z=\widehat V\Psi,\quad
\widehat U=G^{-1}UG,\quad
\widehat V=G^{-1}VG+\dfrac{\ii\Omega}{2}\sg.
\end{equation}
其中 $\widehat U=-\ii\lambda\sg+A(u)$, 且背景矩阵满足
\begin{equation}\label{eq:Rbg}
R_\infty=
\begin{pmatrix}\eta_0&\dfrac{\ii\eta_0Q}{a}\\[6pt]-\dfrac{\ii\eta_0Q}{a}&-\eta_0\end{pmatrix}
=\dfrac{\ii\eta_0}{a}U_\infty(a).
\end{equation}
\begin{equation}\label{eq:F}
\widehat V_\infty=-F(\lambda)U_\infty,\qquad
F(\lambda)=4\beta\lambda^2-\alpha\lambda-2\beta Q^2+
\dfrac{\eta_0}{a(\lambda-a)}.
\end{equation}
\end{proposition}
\begin{proof}
由 $\Phi_z=G_z\Psi+G\Psi_z=VG\Psi$, 时间矩阵的变换具体为
\[
\Psi_z=(G^{-1}VG-G^{-1}G_z)\Psi,\qquad
G^{-1}G_z=-\dfrac{\ii\Omega}{2}\sg.
\]
$G$ 与 $x$ 无关, 因此空间矩阵没有相应的导数项. 对非对角势,
\[
G^{-1}A(q)G=
\begin{pmatrix}0&e^{\ii\Omega z}q\\-e^{-\ii\Omega z}\conj q&0\end{pmatrix}
=A(u),
\]
同一共轭将 $p$ 替换为 $v$, 保持 $\eta$ 不变. 由 $u_z=e^{\ii\Omega z}q_z+\ii\Omega u$, 新光场方程为
\[
u_z=\ii\alpha\left(\dfrac{1}{2}u_{xx}+|u|^2u\right)
+\beta(u_{xxx}+6|u|^2u_x)+2v+\ii\Omega u.
\]

在常背景上, 式 \eqref{eq:Rbg} 由矩阵乘法得到. 时间矩阵的两个独立分量为
\[
(\widehat V_\infty)_{12}
=Q\left(-4\beta\lambda^2+\alpha\lambda+2\beta Q^2
-\dfrac{\eta_0}{a(\lambda-a)}\right).
\]
以及
\begin{align*}
(\widehat V_\infty)_{11}
&=\ii\left(4\beta\lambda^3-\alpha\lambda^2-2\beta Q^2\lambda
+\dfrac{\alpha Q^2}{2}+\dfrac{\eta_0}{\lambda-a}
+\dfrac{\Omega}{2}\right)\\
&=\ii\lambda\left(4\beta\lambda^2-\alpha\lambda-2\beta Q^2+
\dfrac{\eta_0}{a(\lambda-a)}\right),
\end{align*}
最后一步用了 $\dfrac{\alpha Q^2}{2}+\dfrac{\Omega}{2}=\dfrac{\eta_0}{a}$.
故 $\widehat V_\infty=-F U_\infty$.
\end{proof}
有理项 $\dfrac{\eta_0}{a(\lambda-a)}$ 是背景传播因子中介质响应的贡献, 它将在共振点的散射演化与重构中出现.

\section{单值化与背景谱结构}
引入谱曲面 $\rho^2=\lambda^2+Q^2$ 及单值化坐标
\begin{equation}\label{eq:uniform}
w=\lambda+\rho,\quad
\lambda(w)=\dfrac{1}{2}\left(w-\dfrac{Q^2}{w}\right),\quad
\rho(w)=\dfrac{1}{2}\left(w+\dfrac{Q^2}{w}\right).
\end{equation}
第一叶的无穷远由 $\rho\sim\lambda$ 选定, 对应 $w=\infty$; 另一无穷远对应 $w=0$.

\begin{proposition}[背景谱与解析区域]\label{prop:curve}\label{prop:domains}
背景矩阵满足 $U_\infty^2=-(\lambda^2+Q^2)I$. 其振荡谱轮廓、单值化后的轮廓及解析区域分别为
\begin{equation}\label{eq:lambda-spectrum}
\Sigma_\lambda=\mathbb R\cup\ii\lbrack-Q,Q\rbrack.
\end{equation}
\begin{align}
\Sigma&=\mathbb R\cup\{w:|w|=Q\},\label{eq:Sigma}\\
D_\pm&=\left\{w:\ \pm\operatorname{Im}w
\left(1-\dfrac{Q^2}{|w|^2}\right)>0\right\}.\label{eq:domains}
\end{align}
其中 $D_\pm=\{w:\pm\operatorname{Im}\rho(w)>0\}$. 换叶映射为
\begin{equation}\label{eq:sheet}
h(w)=-\dfrac{Q^2}{w},\qquad
\lambda(h(w))=\lambda(w),\quad \rho(h(w))=-\rho(w).
\end{equation}
取背景特征向量矩阵
\begin{equation}\label{eq:E}
E(w)=I-\dfrac{\ii Q}{w}\sigma_1
=\begin{pmatrix}1&-\dfrac{\ii Q}{w}\\[6pt]-\dfrac{\ii Q}{w}&1\end{pmatrix},
\quad \sigma_1=\begin{pmatrix}0&1\\1&0\end{pmatrix}.
\end{equation}
则
\begin{equation}\label{eq:diagE}
U_\infty E=-\ii\rho E\sg,\quad
\gamma(w):=\det E=1+\dfrac{Q^2}{w^2},\quad
E^{-1}=\dfrac{I+\dfrac{\ii Q}{w}\sigma_1}{\gamma(w)}.
\end{equation}
在 $w\ne0,\pm\ii Q,w_a^\pm$ 上, 背景同时 Lax 解为
\begin{equation}\label{eq:theta}
\Psi_0=E(w)e^{-\ii\theta(w;x,z)\sg},\qquad
\theta=\rho(w)[x-F(\lambda(w))z].
\end{equation}
这里共振原像为
\begin{equation}\label{eq:resonance-w}
w_a^\pm=a\pm\sqrt{a^2+Q^2}.
\end{equation}
\end{proposition}
\begin{proof}
矩阵平方的非对角元相消, 对角元均为 $-\lambda^2-Q^2$. 特征值为 $\pm\ii\rho$, 其中 $\rho^2=\lambda^2+Q^2$. 两指数模式无指数增长要求 $\rho\in\mathbb R$. 写 $\lambda=\xi+\ii\nu$, 则
$\rho^2=\xi^2-\nu^2+Q^2+2\ii\xi\nu$. 条件为 $\xi\nu=0$ 及 $\xi^2-\nu^2+Q^2\ge0$, 给出实轴和 $|\nu|\le Q$ 的虚线段. 在端点 $\rho=0$, 矩阵一般为非零幂零矩阵, 基本解包含线性项, 此时不存在两个线性无关的振荡特征向量.

由 $(\rho-\lambda)(\rho+\lambda)=Q^2$ 得到 \eqref{eq:uniform}, 交换谱根符号即得 \eqref{eq:sheet}. 写 $w=re^{\ii\varphi}$, 则
\[
\operatorname{Im}\rho=\dfrac{1}{2}\left(r-\dfrac{Q^2}{r}\right)\sin\varphi,
\]
其零集为实轴与圆 $|w|=Q$, 正负号给出 $D_\pm$.

对特征值 $-\ii\rho$, 将特征向量第一分量取为 1, 第二分量为
$-\dfrac{\ii(\rho-\lambda)}{Q}=-\dfrac{\ii Q}{w}$; 对 $+\ii\rho$ 则交换两分量, 得到 $E$. 其行列式和逆矩阵由二阶矩阵公式给出. 再由 $\theta_x=\rho$、$\theta_z=-\rho F$,
\[
(\Psi_0)_x=U_\infty\Psi_0,\qquad
(\Psi_0)_z=-F U_\infty\Psi_0=\widehat V_\infty\Psi_0.
\]
最后, 解 $\lambda(w)=a$ 得到两个实根 $w_a^\pm$, 且 $w_a^+w_a^-=-Q^2$.
\end{proof}

轮廓各弧以 $D_+$ 位于左侧为正向: 实轴外段朝右, 内段朝左; 上半圆顺时针, 下半圆逆时针. 外部上半平面与内部下半圆盘同属 $D_+$.

后续构造涉及三类有限特殊点. $\pm\ii Q$ 是 $\gamma$ 的零点, 背景特征向量在此退化; $\pm Q$ 是轮廓交点; $w_a^\pm$ 是传播矩阵的共振原像. 空间谱问题在 $w_a^\pm$ 正则, 但当 $z\ne0$ 时, 背景传播指数在这些点具有奇异振荡. 因此共振处采用固定 $z$ 的空间归一化.
